\chapter{Common up--down Krajewski loop: derived cross-word и CP-even fork}

Предыдущая глава показала, что invariant
\(W=\Tr(P_-H_uH_d)\) классифицирует relative up/down geometry, но отсутствует
в direct sum двух square actions. Теперь построим один общий finite graph,
где этот cross-word возникает непосредственно из spectral trace.

\section{Минимальный \texorpdfstring{\(J\)}{J}-парный rectangle}

Первичный rectangle состоит из бимодулей
\begin{equation}
 A=(L_0,R_0),\quad B=(L_0,R_1),\quad
 C=(L_1,R_1),\quad D=(L_1,R_0).
\end{equation}
Edges \(AB,BC,CD,DA\) поочерёдно сохраняют left или right label, поэтому все
проходят order-one condition. Чтобы \(J\) не отождествлял независимые
up/down edges, добавляется транспонированный conjugate rectangle. Полный
оператор имеет восемь graph nodes и dimension \(8\times3=24\) на family
blocks.

На первичном loop фиксируем family factors
\begin{equation}
 AB:P_-,
 \qquad
 BC:H_u,
 \qquad
 CD:H_d,
 \qquad
 DA:e^{i\theta}I_3.
\end{equation}
Противоположные matrix elements задаются self-adjointness, а второй rectangle
--- реальной структурой. Прямой matrix test подтверждает одновременно
self-adjointness, \(JD_F=D_FJ\), odd grading и order-one edge criterion.

\section{Cross-word из одного trace}

Для каждой из трёх relative orbit полный quartic trace имеет универсальный
вид
\begin{equation}
 \Tr D_F^4
 =
 104+16\cos\theta\,\Tr(P_-H_uH_d).
\end{equation}
Следовательно, mixed coefficient не вводится вручную: он является
кратностью ориентированного четырёхрёберного loop и имеет фиксированный знак.

Подстановка \(W=-1/2,0,+1/2\) даёт
\begin{align}
 W=-\frac12:&\quad \Tr D_F^4=104-8\cos\theta,\\
 W=0:&\quad \Tr D_F^4=104,\\
 W=+\frac12:&\quad \Tr D_F^4=104+8\cos\theta.
\end{align}
При положительном quartic coefficient minima равны
\begin{equation}
 (W,\theta)=
 \left(-\frac12,0\right)
 \quad\text{или}\quad
 \left(+\frac12,\pi\right),
\end{equation}
и имеют одинаковое значение \(96\). Commuting class \(W=0\) остаётся на
уровне \(104\) и исключается.

\begin{theorem}[Derived noncommuting-orbit selection]
Один общий \(J\)-парный Krajewski rectangle автоматически порождает
cross-word и спектрально исключает commuting relative orbit. Тем самым три
физические up/down architecture сокращаются до двух noncommuting branches
без введения нового \(\kappa_{ud}\).
\end{theorem}

\section{Почему CP ещё не закрыта}

Loop phase релаксирует вместе с orbit sign. В двух minima она равна \(0\)
или \(\pi\), поэтому оба loop products вещественны. Transformation
\begin{equation}
 (W,\theta)\longmapsto(-W,\theta+\pi)
\end{equation}
сохраняет minimum energy. Следовательно, action выбирает noncommutativity, но
не orientation и не family-nonfactorizable CP phase.

\begin{nogocriterion}[CP-even cross-loop]
Наличие derived cross-word не разрешает объявить CKM phase выведенной.
Нужен дополнительный orientation-sensitive datum, который запрещает
компенсацию sign \(W\) сдвигом \(\theta\mapsto\theta+\pi\), причём этот datum
должен возникать из того же общего operator, а не из внешнего complex
coefficient.
\end{nogocriterion}

Гейт является реальным продвижением: mixed up/down interaction впервые
получен из одного \(D_F\), а commuting branch устранена. Открытым остаётся
двухветвевой CP-even fork. Exact \(24\times24\) ledger сохранён в
\path{s2t_v4_common_updown_krajewski_loop_gate_results.json}.